% ============================================================
%  HPSC 2026 – Assignment 1 White Paper
%  IEEE two-column conference style (article + geometry)
% ============================================================
\documentclass[10pt,twocolumn]{article}

\usepackage[a4paper,top=19mm,bottom=43mm,left=15.5mm,right=15.5mm,columnsep=4.5mm]{geometry}
\usepackage{amsmath,amssymb}
\usepackage{graphicx}
\graphicspath{{figures/}}
\usepackage{booktabs}
\usepackage[hidelinks]{hyperref}
\usepackage[table]{xcolor}
\usepackage{listings}
\usepackage{caption}
\usepackage{subcaption}
\usepackage{algorithm}
\usepackage{algpseudocode}
\usepackage{siunitx}
\usepackage{cite}
\usepackage[T1]{fontenc}
\usepackage{microtype}
\usepackage{times}
\usepackage{titlesec}
\usepackage{abstract}
\usepackage{fancyhdr}
\usepackage{setspace}

\titleformat{\section}{\normalfont\bfseries}{\Roman{section}.}{0.5em}{\MakeUppercase}
\titleformat{\subsection}{\normalfont\itshape}{\Alph{subsection}.}{0.5em}{}
\titlespacing{\section}{0pt}{8pt plus2pt minus1pt}{4pt}
\titlespacing{\subsection}{0pt}{6pt plus1pt}{2pt}

\renewcommand{\abstractnamefont}{\normalfont\bfseries}
\renewcommand{\abstracttextfont}{\normalfont\small}
\setlength{\absleftindent}{0pt}
\setlength{\absrightindent}{0pt}

\captionsetup{font=small,labelfont=bf,skip=4pt}

\lstset{
  language=C++,basicstyle=\ttfamily\scriptsize,
  keywordstyle=\color{blue!75!black}\bfseries,
  commentstyle=\color{green!50!black},stringstyle=\color{orange!80!black},
  numbers=left,numberstyle=\tiny\color{gray},numbersep=4pt,
  breaklines=true,frame=single,xleftmargin=8pt,framexleftmargin=8pt,
  aboveskip=4pt,belowskip=4pt
}

\definecolor{rowgray}{gray}{0.93}

\newcommand{\vv}[1]{\boldsymbol{#1}}
\newcommand{\norm}[1]{\left\|#1\right\|}

\pagestyle{fancy}\fancyhf{}
\renewcommand{\headrulewidth}{0pt}
\fancyfoot[C]{\small\thepage}

\setlength{\columnseprule}{0.4pt}
\def\columnseprulecolor{\color{gray!40}}

\begin{document}

\twocolumn[{%
\begin{@twocolumnfalse}
  \begin{center}
    {\large\bfseries OpenMP Parallelisation of a Three-Dimensional\\[2pt]
      Discrete Element Method Particle Solver}\\[8pt]
    {\normalsize\itshape Tushar}\\[2pt]
    {\small School of Mechanical and Material Sciences, IIT Mandi \quad \texttt{Tusharkharod366@gmail.com}}\\[10pt]
  \end{center}
  \begin{abstract}
  \noindent
  We present the implementation, verification, and OpenMP parallelisation
  of a three-dimensional Discrete Element Method (DEM) solver for spherical
  particles confined in a cuboidal domain.
  The solver uses a linear spring-dashpot contact model and semi-implicit
  Euler time integration.
  Three verification tests confirm numerical correctness, and a convergence
  study confirms first-order accuracy (measured slope $\approx1.00$).
  A cell-linked list neighbour search reduces contact-detection cost from
  $\mathcal{O}(N^2)$ to $\mathcal{O}(N)$ average complexity.
  Speedup and parallel efficiency are measured for
  $N\in\{200,1000,5000\}$ particles on up to $p=8$ threads, achieving a
  peak speedup of $S\approx6.22$ (efficiency $E\approx77.7\%$) at $N=5000$.
  Combined, algorithmic and hardware-level optimisations reduce total runtime
  by approximately two orders of magnitude for large-$N$ systems.
  \end{abstract}
  \vspace{10pt}
\end{@twocolumnfalse}
}]

%─────────────────────────────────────────────────────────────
\section{Introduction}\label{sec:intro}

Discrete Element Methods (DEM) model granular materials by explicitly
tracking every particle and computing pairwise contact forces~\cite{cundall1979}.
Originally introduced by Cundall and Strack, DEM is now indispensable in
process engineering, geomechanics, and pharmaceutical
manufacturing~\cite{dBonn2017granular}.
The dominant bottleneck is the $\mathcal{O}(N^2)$ contact-detection loop,
motivating both algorithmic (neighbour search) and hardware-level
(OpenMP) optimisations.

Section~\ref{sec:model} presents the physical model;
\ref{sec:algo} the numerical algorithm;
\ref{sec:verify} verification;
\ref{sec:profile} profiling;
\ref{sec:parallel} OpenMP parallelisation;
\ref{sec:perf} performance analysis.

%─────────────────────────────────────────────────────────────
\section{Mathematical Formulation}\label{sec:model}

\subsection{Equations of Motion}

Consider $N$ spherical particles with positions $\vv{x}_i$, velocities
$\vv{v}_i$, masses $m_i$, and radii $R_i$.
Newton's second law gives
\begin{align}
  m_i \frac{d\vv{v}_i}{dt} &= \vv{F}_i ,\label{eq:newton}\\
  \frac{d\vv{x}_i}{dt}     &= \vv{v}_i ,\label{eq:kin}
\end{align}
where
\begin{equation}
  \vv{F}_i = m_i\vv{g}
           + \sum_{j\neq i}\vv{F}^c_{ij}
           + \vv{F}^w_i .\label{eq:force_total}
\end{equation}

\subsection{Spring-Dashpot Contact Model}

The normal overlap between particles $i$ and $j$ is
\begin{equation}
  \delta_{ij} = R_i+R_j-\norm{\vv{r}_{ij}},\quad
  \vv{r}_{ij}=\vv{x}_j-\vv{x}_i .\label{eq:overlap}
\end{equation}
Contact exists when $\delta_{ij}>0$.  The force is
\begin{equation}
  \vv{F}^c_{ij}=F_n\hat{\vv{n}}_{ij},\quad
  F_n=\max\!\bigl(0,\;k_n\delta_{ij}-\gamma_n v_{n,ij}\bigr),
  \label{eq:contact}
\end{equation}
where $\hat{\vv{n}}_{ij}=\vv{r}_{ij}/\norm{\vv{r}_{ij}}$,
$v_{n,ij}=(\vv{v}_j-\vv{v}_i)\cdot\hat{\vv{n}}_{ij}$,
$k_n$ is normal stiffness, and $\gamma_n$ is damping.
The clamp prevents spurious attractive forces.

\subsection{Wall Contact}

Domain: $[0,L_x]\times[0,L_y]\times[0,L_z]$.
For the lower $z$ wall, overlap $\delta^w=R_i-z_i$ gives
$F^w_z=\max(0,k_n\delta^w-\gamma_n(-v_{z,i}))$.
Analogous expressions apply to all six faces.

%─────────────────────────────────────────────────────────────
\section{Numerical Algorithm}\label{sec:algo}

\subsection{Time Integration}

Semi-implicit Euler (first-order):
\begin{align}
  \vv{v}^{n+1}_i &= \vv{v}^n_i + \frac{\vv{F}^n_i}{m_i}\Delta t,
  \label{eq:vel_update}\\
  \vv{x}^{n+1}_i &= \vv{x}^n_i + \vv{v}^{n+1}_i\Delta t.
  \label{eq:pos_update}
\end{align}
Using $\vv{v}^{n+1}$ in the position update improves stability over
explicit Euler at no additional cost.

\subsection{Timestep Algorithm}

\begin{algorithm}[H]
\small
\caption{Single DEM Timestep}\label{alg:timestep}
\begin{algorithmic}[1]
  \ForAll{$i$}
    \State $\vv{F}_i\!\leftarrow\!\vv{0}$;
           $\;\vv{F}_i\!\mathrel{+}=\!m_i\vv{g}$
  \EndFor
  \ForAll{pairs $(i,j)$, $i<j$}
    \State Compute $\delta_{ij}$ [Eq.~\eqref{eq:overlap}]
    \If{$\delta_{ij}>0$}
      \State $\vv{F}_i\mathrel{+}=\vv{F}^c_{ij}$;
             $\;\vv{F}_j\mathrel{-}=\vv{F}^c_{ij}$
    \EndIf
  \EndFor
  \ForAll{$i$, walls $w$}
    \State $\vv{F}_i\mathrel{+}=\vv{F}^w$
  \EndFor
  \ForAll{$i$}
    \State Update $\vv{v}_i$, $\vv{x}_i$
           [Eqs.~\eqref{eq:vel_update}--\eqref{eq:pos_update}]
  \EndFor
\end{algorithmic}
\end{algorithm}

\subsection{Cell-Linked List Neighbour Search}

Brute force evaluates $N(N{-}1)/2$ pairs per step ($\mathcal{O}(N^2)$).
The cell-linked list~\cite{hockney1988} divides the domain into cells of
side $\ell_c=2R_{\max}$ and restricts checks to the $3^3=27$ neighbouring
cells, yielding $\mathcal{O}(N)$ average complexity.

%─────────────────────────────────────────────────────────────
\section{Verification and Testing}\label{sec:verify}

\subsection{Test 1: Free Fall}

A single particle ($m=\SI{1}{kg}$, $R=\SI{0.05}{m}$) released from rest at
$z_0=\SI{5}{m}$.
Analytical: $z(t)=z_0-\tfrac{1}{2}gt^2$.
Figure~\ref{fig:freefall} shows excellent agreement over $t\in[0,1]\,\si{s}$.

\begin{figure}[t]
  \centering
  \includegraphics[width=\linewidth]{free_fall_trajectory.pdf}
  \caption{Free-fall: numerical vs.\ analytical $z(t)$ (top)
           and $v_z(t)$ (bottom). $\Delta t=10^{-4}\,\si{s}$.}
  \label{fig:freefall}
\end{figure}

\subsection{Test 2: Constant Velocity}

With $\vv{g}=\vv{0}$ and initial $\vv{v}_0=(1,\,0.5,\,0.3)\,\si{m\,s^{-1}}$,
the maximum positional error was $<10^{-10}\,\si{m}$, confirming no
spurious acceleration.

\subsection{Test 3: Bouncing Particle}

A particle dropped from $z=\SI{2}{m}$ (Fig.~\ref{fig:bounce}).
Maximum wall penetration $<\!1\%$ of $R$.
Rebound height decreases monotonically confirming energy dissipation.

\begin{figure}[t]
  \centering
  \includegraphics[width=\linewidth]{bounce.pdf}
  \caption{Bounce test: height $z(t)$ (top), kinetic energy (bottom).
           $k_n=10^5\,\si{N\,m^{-1}}$,
           $\gamma_n=\SI{50}{N.s.m^{-1}}$.}
  \label{fig:bounce}
\end{figure}

\subsection{Convergence Study}

Figure~\ref{fig:convergence} plots RMSE of $z(t)$ vs $\Delta t$ for seven
step sizes.  Measured slope $\approx1.00$ confirms first-order accuracy.

\begin{figure}[t]
  \centering
  \includegraphics[width=\linewidth]{convergence.pdf}
  \caption{RMSE vs $\Delta t$; measured slope $\approx1.00$.}
  \label{fig:convergence}
\end{figure}

%─────────────────────────────────────────────────────────────
\section{Profiling Results}\label{sec:profile}

Table~\ref{tab:profile} shows the runtime breakdown for $N=1000$ over
$50\,000$ timesteps (single core, \texttt{-O2 -march=native}).
Contact detection accounts for $\approx\!95\%$ of total time.
Enabling \texttt{-O2} gave a $3.8\times$ speedup over the unoptimised
build before any parallelisation.

\begin{table}[t]
\centering
\caption{Serial Runtime Breakdown ($N=1000$, $50\,000$ steps)}\label{tab:profile}
\setlength{\tabcolsep}{5pt}
\renewcommand{\arraystretch}{1.1}
\begin{tabular}{lrr}
\toprule
Routine & Time (s) & \% \\
\midrule
\rowcolor{rowgray}
\texttt{compute\_particle\_contacts} & 7.79 & 95.0 \\
\texttt{compute\_wall\_contacts}     & 0.06 &  0.7 \\
\rowcolor{rowgray}
\texttt{semi\_implicit\_euler}       & 0.02 &  0.2 \\
\texttt{write\_output}               & 0.33 &  4.0 \\
\midrule
Total                                & 8.20 & 100  \\
\bottomrule
\end{tabular}
\end{table}

%─────────────────────────────────────────────────────────────
\section{Parallel Implementation}\label{sec:parallel}

\subsection{Thread-Local Force Buffers}

The contact loop writes to $\vv{F}_i$ and $\vv{F}_j$ per pair, creating
a race condition under plain \texttt{\#pragma omp for}.
Two strategies were evaluated:

\begin{enumerate}\setlength\itemsep{0pt}
  \item \textit{Atomic updates}---simple but serialises memory traffic.
  \item \textit{Thread-local buffers}---each thread writes to a private
        size-$N$ array; arrays are reduced serially after the parallel
        region.  Lock-free on the hot path; superior for large $N$.
\end{enumerate}

Strategy~(2) is adopted with \texttt{schedule(dynamic,32)}.

\begin{lstlisting}[caption={OpenMP contact loop (simplified).},label={lst:omp}]
vector<vector<Vec3>> tf(NTHREADS,
    vector<Vec3>(N, Vec3(0,0,0)));
#pragma omp parallel
{
  int tid = omp_get_thread_num();
  #pragma omp for schedule(dynamic,32) nowait
  for (int i = 0; i < N-1; ++i)
    for (int j = i+1; j < N; ++j) {
      // ... compute contact force Fc ...
      tf[tid][i] += Fc;
      tf[tid][j] -= Fc;
    }
} // barrier
for (int t=0; t<NTHREADS; ++t)
  for (int i=0; i<N; ++i)
    particles[i].force += tf[t][i];
\end{lstlisting}

Correctness confirmed: max relative KE difference between serial and
OpenMP runs was $<10^{-12}$.

%─────────────────────────────────────────────────────────────
\section{Performance Analysis}\label{sec:perf}

\subsection{Scaling Study}

\begin{equation}
  S(p)=\frac{T_1}{T_p},\qquad E(p)=\frac{S(p)}{p}.
\end{equation}

Table~\ref{tab:scaling} lists measured $T_p$ and $S(p)$;
Figure~\ref{fig:scaling} shows the curves.

\begin{table}[t]
\centering
\caption{$T_p$ (s) and $S(p)$ for varying $N$ and $p$.}\label{tab:scaling}
\setlength{\tabcolsep}{3.5pt}
\renewcommand{\arraystretch}{1.1}
\begin{tabular}{crrrrrr}
\toprule
 & \multicolumn{2}{c}{$N=200$}
 & \multicolumn{2}{c}{$N=1000$}
 & \multicolumn{2}{c}{$N=5000$}\\
\cmidrule(lr){2-3}\cmidrule(lr){4-5}\cmidrule(lr){6-7}
$p$ & $T_p$ & $S$ & $T_p$ & $S$ & $T_p$ & $S$\\
\midrule
\rowcolor{rowgray}
1 & 0.38 & 1.00 & 8.20 & 1.00 & 105.0 & 1.00\\
2 & 0.25 & 1.54 & 4.59 & 1.79 &  53.6 & 1.96\\
\rowcolor{rowgray}
4 & 0.18 & 2.18 & 2.79 & 2.94 &  30.0 & 3.50\\
8 & 0.14 & 2.74 & 1.99 & 4.12 &  16.9 & 6.22\\
\bottomrule
\end{tabular}
\end{table}

\begin{figure}[t]
  \centering
  \includegraphics[width=\linewidth]{scaling_speedup_efficiency.pdf}
  \caption{Left: speedup $S(p)$. Right: efficiency $E(p)\,(\%)$.
           Dotted: ideal scaling.}
  \label{fig:scaling}
\end{figure}

At $N=5000$, $p=8$: $S\approx6.22$, $E\approx77.7\%$.
Sub-linear scaling at $N=200$ ($S_{\max}=2.74$) reflects a large serial
fraction ($f\approx28\%$, Amdahl limit $1/f\approx3.6$).
For $N=5000$, $f\approx4\%$ permits much higher speedup.

\subsection{Kinetic Energy Evolution}

Figure~\ref{fig:ke} shows total kinetic energy for settling simulations.
The initial peak corresponds to free-fall acceleration; the subsequent
monotonic decay reflects dashpot dissipation.

\begin{figure}[t]
  \centering
  \includegraphics[width=\linewidth]{kinetic_energy.pdf}
  \caption{Kinetic energy vs time, $N\in\{200,1000,5000\}$.}
  \label{fig:ke}
\end{figure}

\subsection{Particle Configuration}

Figure~\ref{fig:snap} shows a snapshot at $t=0.8\,\si{s}$ ($N=500$),
coloured by speed.  Fast particles cluster near the floor;
slower particles are settled at greater heights.

\begin{figure}[t]
  \centering
  \includegraphics[width=0.88\linewidth]{snapshot.pdf}
  \caption{Particle snapshot at $t=0.8\,\si{s}$, $N=500$.
           Colour = speed (m\,s$^{-1}$).}
  \label{fig:snap}
\end{figure}

\subsection{Neighbour Search}

Figure~\ref{fig:nbr} compares per-timestep cost.
The crossover is at $N\approx850$; above this the neighbour search gives
$\approx12\times$ speedup at $N=10\,000$.

\begin{figure}[t]
  \centering
  \includegraphics[width=\linewidth]{neighbor_runtime.pdf}
  \caption{Runtime: $\mathcal{O}(N^2)$ vs neighbour search.
           Crossover at $N\approx850$.}
  \label{fig:nbr}
\end{figure}

%─────────────────────────────────────────────────────────────
\section{Discussion and Conclusions}\label{sec:conclusions}

A modular 3-D DEM solver was implemented in C++ and verified against
analytical solutions for free fall, constant velocity, and wall bounce.
The semi-implicit Euler integrator shows first-order convergence
(slope~$\approx1.00$).
OpenMP parallelisation with thread-local force buffers achieves
$S\approx6.22$ ($E\approx77.7\%$) at $N=5000$, $p=8$.
The cell-linked list provides $\approx12\times$ speedup at $N=10\,000$,
with crossover at $N\approx850$.
Combined, optimisations reduce total runtime by $\sim\!2$ orders of
magnitude for large systems.

Future directions include rotational dynamics, non-spherical particles,
and MPI parallelism for $N>10^6$.

\section*{Acknowledgements}
LLM assistance (Claude, Anthropic) used for scaffolding; all results
independently verified.
Code: \url{https://github.com/TuShhar0007/dem_hpsc2026}

\bibliographystyle{IEEEtran}
\begin{thebibliography}{9}

\bibitem{cundall1979}
P.~A.~Cundall and O.~D.~L.~Strack,
``A discrete numerical model for granular assemblies,''
\textit{G\'eotechnique}, vol.~29, no.~1, pp.~47--65, 1979.

\bibitem{dBonn2017granular}
D.~Bonn \textit{et al.},
``Yield stress materials in soft condensed matter,''
\textit{Rev.\ Mod.\ Phys.}, vol.~89, p.~035005, 2017.

\bibitem{hockney1988}
R.~W.~Hockney and J.~W.~Eastwood,
\textit{Computer Simulation Using Particles}.
IOP Publishing, 1988.

\bibitem{omp2018}
OpenMP Architecture Review Board,
\textit{OpenMP Application Programming Interface, v5.0}, 2018.

\end{thebibliography}

\end{document}
